\thispagestyle{empty}

\begin{center}
  {\Large\bfseries
    ChatGPT's Review of the Manuscript
  }
\end{center}

\begin{center}
\small
\textit{
The following review was generated by ChatGPT after reading the manuscript.
}
\end{center}

\begin{quotation}\small
At first sight, \textit{Quantum Resonance as a Global Connection Problem} looks like an unusually broad monograph. In roughly 450 pages it moves from operator domains, spectral measures, and scattering theory to exact WKB, complex scaling, few-body nuclear reactions, exceptional points, and black-hole quasinormal modes.

Yet its breadth is not what makes the book unusual.

\begin{center}
\textbf{The book is strange not because it is sprawling, but because it follows one organizing principle far beyond the point where most textbooks would have changed languages.}
\end{center}

That principle is that a resonance should not be defined first as a complex eigenvalue of a non-self-adjoint Hamiltonian. One begins instead with a self-adjoint quantum problem and its scattering boundary data, constructs a global connection coefficient, and analytically continues it. The resonance condition $\Cincoming(E_R)=0$ comes first. Resolvent and scattering poles, Gamow states, Zel'dovich regularization, complex-scaled eigenvectors, and rigged-Hilbert-space functionals then appear as different realizations of the same continued analytic datum.

This reversal of the usual order is the real subject of the book. It also explains the unusually substantial mathematical opening. Domains, operator topologies, spectral measures, direct integrals, and \vN\ algebras are not included merely for completeness; they specify what remains meaningful when representations, continuum limits, and finite-dimensional approximations are changed.

A particularly effective expression of this philosophy is the ``truncation contract.'' Since every actual computation is finite, the important question is not whether one truncates, but under what statement the finite problem represents the intended infinite-dimensional one: which parent operator, which projection or basis, which topology of convergence, which observable, and which order of limits.

\begin{center}
\textbf{In this book, ``the matrix has a stable complex eigenvalue'' is never allowed to become a substitute for ``the physical resonance has been identified.''}
\end{center}

Exact WKB provides the main constructive machinery. Local Riccati solutions are Borel resummed, transported through Stokes sectors, and assembled into a global connection problem. Scattering theory identifies the physically relevant coefficient, while complex scaling and related methods provide alternative representations of the resulting pole.

One qualification is important. The \textit{connection-first} viewpoint is more universal than exact WKB itself. For one-dimensional and radial differential equations, exact WKB genuinely acts as a constructive spine. In coupled-channel nuclear reactions, many-body problems, and possible extensions toward quantum field theory, the deeper common structure is instead the broader one of asymptotic sectors, projections, boundary conditions, resolvents, and analytic continuation.

The book is commendably explicit about this limitation. Black-hole complex scaling still lacks a general theorem comparable to the standard analytic-dilation setting; threshold problems require uniform treatment when poles and branch points collide; exact-WKB residues are less developed than pole positions; continuum exceptional points require a regulator-independent formulation; and ``monodromy renormalization'' is presented as an organizing principle rather than a completed universal theory.

This restraint is important because the unifying ambition is otherwise very strong.

The later chapters test that ambition in increasingly different settings. Radial singularities add Frobenius and monodromy data; few-body reactions replace a single channel by projected channel spaces and continuum discretization; exceptional points turn a simple resonance zero into a multiple one and hence a Jordan structure; black-hole horizons replace ordinary asymptotic scattering by horizon radiation conditions. The notation and techniques change, but the question remains the same: what is the global analytic datum, and which parts of a calculation are invariant under a change of representation?

The continuum is equally central. The manuscript repeatedly refuses to treat it as disposable background after a pole has been found. Branch cuts determine sheets, rotated continua complete complex-scaled representations, threshold structure controls pole approximations, and continuum response contributes to observables.

The price of this architecture is density. The manuscript is formally self-contained from modest prerequisites, but a reader meeting functional analysis, resurgence, few-body reaction theory, and black-hole perturbation theory for the first time will face a steep conceptual gradient. This difficulty, however, comes less from excess material than from compression.

Concepts are introduced once in canonical chapters and reused later; repeated introductions are removed; research calculations are dismantled and rebuilt in prerequisite order.

\begin{center}
\textbf{The remarkable compression of this book is not that it omits mathematics, but that it refuses to restart the theory every time the physical system or representation changes.}
\end{center}

For this reason, I would regard the work as a graduate textbook with genuine monographic coherence rather than as a review or encyclopedia. Most of its ingredients are established knowledge. Its originality lies primarily in deciding what should be regarded as primitive, what is representation-dependent, and in what order the subject should be reconstructed.

Its central claim can therefore be stated very simply:

\begin{center}
\textbf{A resonance is not fundamentally a complex eigenvalue. It is a global analytic connection datum, and its legitimate realizations must agree on the same underlying resonance information.}
\end{center}

Whether this viewpoint becomes standard is impossible to predict. But the manuscript does succeed in making the alternative viewpoint---treating scattering poles, Gamow states, complex scaling, and non-self-adjoint eigenvalues as essentially independent notions---look increasingly unnatural.

That is the real achievement of the book.
\end{quotation}

\vspace{1em}

\begin{flushright}
--- ChatGPT
\end{flushright}

\vfill

\begin{flushright}
\textbf{Reply to ChatGPT}

\textit{What would remain if modern theoretical physics were compressed\\
to a Landau minimum?}

--- Okuto Morikawa
\end{flushright}